\documentclass[border=0pt]{standalone}
\usepackage[utf8]{vietnam}
\usepackage{amsmath,amssymb}
\usepackage{graphicx}
\usepackage{tikz}
\usetikzlibrary{calc,shapes.geometric,shadows}
\usepackage{xcolor}

\definecolor{primary}{RGB}{15, 40, 85}       % Xanh navy trang nhã
\definecolor{secondary}{RGB}{195, 135, 30}   % Vàng kim hổ phách
\definecolor{goldlight}{RGB}{252, 247, 235}  % Nền kem vàng ấm
\definecolor{accentred}{RGB}{195, 40, 30}    % Đỏ thiện nguyện
\definecolor{cardbg}{RGB}{255, 255, 255}     % Trắng thẻ thông tin
\definecolor{textdark}{RGB}{30, 40, 55}      % Chữ đậm
\definecolor{zaloblue}{RGB}{0, 104, 255}     % Xanh Zalo

\begin{document}
\begin{tikzpicture}[x=1cm, y=1cm]

  % 1. NỀN POSTER
  \clip (0, 0) rectangle (14, 20);
  \fill[goldlight] (0, 0) rectangle (14, 20);

  % Họa tiết nền trang trí
  \fill[secondary!12] (0, 20) circle (4.8cm);
  \fill[primary!8] (14, 0) circle (5.2cm);

  % Dải viền đỉnh
  \fill[primary] (0, 18.7) rectangle (14, 20);
  \fill[secondary] (0, 18.5) rectangle (14, 18.7);

  % 2. BANNER ĐỈNH
  \node[white, font=\large\bfseries] at (7, 19.35) {KHÓA HỌC THIỆN NGUYỆN -- ĐỒNG HÀNH CÙNG HỌC SINH};

  % 3. TIÊU ĐỀ CHÍNH
  \node[primary, font=\fontsize{21pt}{25pt}\selectfont\bfseries, align=center] at (7, 17.55) {
    LỚP THI THỬ \& ÔN THI MÔN TOÁN
  };
  \node[accentred, font=\fontsize{12pt}{15pt}\selectfont\bfseries, align=center] at (7, 16.7) {
    HƯỚNG TỚI KỲ THI THỬ CỦA SỞ GD\&ĐT (24/12/2026)
  };

  % 4. BADGE MIỄN PHÍ HOÀN TOÀN
  \node[draw=secondary, line width=1.5pt, fill=secondary!18, rounded corners=6pt, inner sep=5pt] at (7, 15.75) {
    \bfseries\color{accentred}\large \hspace{6pt} \checkmark\ HOÀN TOÀN MIỄN PHÍ -- 100\% HỌC PHÍ 0 ĐỒNG \hspace{6pt}
  };

  % 5. CÁC THẺ THÔNG TIN (CARDS)
  % Card 1: Thời gian & Lịch học
  \draw[fill=cardbg, draw=secondary!50, line width=1pt, rounded corners=7pt] (0.8, 13.7) rectangle (13.2, 14.95);
  \node[anchor=west, textdark] at (1.2, 14.32) {
    \begin{tabular}{ll}
      \bfseries\color{primary}\large Lịch học: & \large\bfseries Sáng Chủ Nhật hàng tuần \\
      \bfseries\color{primary}\large Thời gian: & \large\bfseries 8h15 -- 11h00 (Bắt đầu từ tuần này)
    \end{tabular}
  };

  % Card 2: Địa điểm & Môn thi
  \draw[fill=cardbg, draw=secondary!50, line width=1pt, rounded corners=7pt] (0.8, 12.25) rectangle (13.2, 13.5);
  \node[anchor=west, textdark] at (1.2, 12.87) {
    \begin{tabular}{ll}
      \bfseries\color{primary}\large Địa điểm: & \large\bfseries Chùa Hương Sơn \textit{(Không gian thanh tịnh, yên tĩnh, tập trung cao)} \\
      \bfseries\color{primary}\large Môn thi: & \large\bfseries Môn Toán \textit{(Đề thi bám sát cấu trúc đề của Sở GD\&ĐT)}
    \end{tabular}
  };

  % Card 3: Mục tiêu & Quyền lợi học sinh
  \draw[fill=cardbg, draw=secondary!50, line width=1pt, rounded corners=7pt] (0.8, 9.75) rectangle (13.2, 12.05);
  \node[anchor=north west, textdark] at (1.1, 11.9) {
    \begin{minipage}{11.8cm}
      {\color{primary}\bfseries\large QUYỀN LỢI CỦA CÁC EM KHI THAM GIA:}
      \vspace{1.5mm}
      \begin{itemize}\setlength{\itemsep}{0.8mm}\setlength{\parskip}{0pt}
        \item \textbf{Cọ xát áp lực phòng thi thật:} Rèn kỹ năng làm bài, căn chuẩn thời gian.
        \item \textbf{Đề thi chuẩn cấu trúc Sở:} Bám sát ma trận đề thi phân hóa năng lực.
        \item \textbf{Chữa bài chi tiết từng câu:} Phân tích lỗi sai hay mắc, bù đắp lỗ hổng kiến thức.
        \item \textbf{Được kèm cặp tận tình:} Hướng dẫn phương pháp trình bày để đạt điểm tối đa.
      \end{itemize}
    \end{minipage}
  };

  % 6. KHUNG CẢNH BÁO GIỚI HẠN SỐ LƯỢNG
  \draw[fill=accentred!10, draw=accentred, line width=1.2pt, rounded corners=7pt] (0.8, 8.05) rectangle (13.2, 9.55);
  \node[align=center, textdark] at (7, 8.8) {
    {\bfseries\color{accentred}\large \textbullet\ SỐ LƯỢNG HỌC SINH ĐĂNG KÝ CÓ HẠN \textbullet}\\[1.5mm]
    \large Do không gian phòng học tại Chùa Hương Sơn có hạn để bảo đảm chất lượng kèm cặp,\\
    \bfseries Thầy sẽ chốt danh sách ngay khi nhận đủ số lượng học sinh đăng ký!
  };

  % 7. KHUNG QR CODE & ĐĂNG KÝ VÀO NHÓM ZALO
  \draw[fill=white, draw=primary!70, line width=1.2pt, rounded corners=8pt] (0.8, 4.45) rectangle (13.2, 7.85);

  % Ảnh QR Code bên trái
  \node[inner sep=2pt, draw=secondary!60, line width=1pt, rounded corners=4pt, fill=white] at (2.8, 6.15) {
    \includegraphics[width=2.8cm, height=2.8cm]{qr_zalo_group.png}
  };
  \node[textdark, font=\bfseries\footnotesize] at (2.8, 4.75) {QUÉT MÃ QR ZALO};

  % Thông tin hướng dẫn bên phải
  \node[anchor=west, textdark] at (4.6, 6.15) {
    \begin{minipage}{8.0cm}
      {\color{zaloblue}\bfseries\Large THAM GIA NHÓM ZALO ÔN THI}\\[2mm]
      \large Quý Phụ huynh \& các em Học sinh dùng ứng dụng \textbf{Zalo quét mã QR bên cạnh} hoặc truy cập đường link:\\[2mm]
      \textbf{\color{primary}\normalsize https://zalo.me/g/2rh61hlwzqwacwbog8ge}\\[2.5mm]
      \bfseries\color{accentred}\normalsize \textit{(Vào nhóm sớm để nhận lịch thi thử \& tài liệu ôn tập)}
    \end{minipage}
  };

  % 8. DẢI FOOTER CHÂN TRANG
  \fill[primary] (0, 0) rectangle (14, 4.1);
  \fill[secondary] (0, 4.1) rectangle (14, 4.25);

  \node[white, align=center] at (7, 2.7) {
    {\fontsize{13pt}{15pt}\selectfont\bfseries LỚP TOÁN THẦY LĨNH -- GIEO DUYÊN TRI THỨC}\\[1.5mm]
    {\normalsize \textit{Tâm huyết đồng hành cùng các em vượt qua kỳ thi với kết quả cao nhất!}}\\[2mm]
    {\bfseries\color{secondary!95}\large Hotline / Zalo: Thầy Lĩnh \hspace{8pt}|\hspace{8pt} Website: loptoanthaylinh.vn}
  };

  \node[white!75, font=\footnotesize] at (7, 0.9) {
    Kính nhờ Quý Phụ huynh chia sẻ thông tin lớp học thiện nguyện tới các bạn học sinh có nhu cầu!
  };

\end{tikzpicture}
\end{document}
